% !TEX root = ../main.tex

%----------------------------------------------------------------------
\subsection{Defect engineering in two-dimensional materials}
%----------------------------------------------------------------------

Point defects---substitutional impurities, interstitials, and
vacancies---fundamentally govern the electronic, optical, and catalytic
properties of two-dimensional (2D) materials~\cite{imp2d2023}.
In transition-metal dichalcogenides (TMDs) such as MoS$_2$ and WSe$_2$,
a single dopant atom can shift the Fermi level, create mid-gap trap states,
or activate catalytic edge sites.
The defect formation energy $\ef$ quantifies the thermodynamic cost of
introducing such a defect and serves as the primary descriptor for
high-throughput defect screening.

However, the combinatorial space of possible host--dopant configurations
is enormous.
The Impurities in 2D Materials Database (IMP2D)~\cite{imp2d2023} covers 44
monolayer hosts doped with 65 elements, yielding a $44\times65 = 2{,}860$
matrix of which only 91.2\% entries have converged DFT results
(10\,641 configurations).
Each DFT supercell calculation requires hours of CPU time, making exhaustive
enumeration impractical for the remaining gaps and for extensions to larger
supercells or additional host families.

%----------------------------------------------------------------------
\subsection{Limitations of existing approaches}
%----------------------------------------------------------------------

Machine-learning (ML) surrogate models offer orders-of-magnitude speedup
over DFT.
The Atomistic Line Graph Neural Network (ALIGNN)~\cite{alignn2021}, a
state-of-the-art crystal property predictor with 4.03\,M parameters, achieves
an MAE of 0.540\eV{} on IMP2D under random 80/10/10 splitting.
El~Alouani \textit{et al.}~\cite{elalouani2024} applied gradient-boosted
trees with hand-crafted Jacobi--Legendre descriptors to the same database,
reaching 0.518\eV{}.
Kazeev \textit{et al.}~\cite{kazeev2023} introduced sparse defect
representations for six host materials, achieving 42\,meV/site (a per-site
metric not directly comparable to per-structure MAE).
In the 3D domain, defect-specific GNNs have been proposed for bulk
semiconductors~\cite{rahman2024,witman2023}, and recent work demonstrates
that charged-defect energies can be predicted from crystal
structure alone~\cite{kiyohara2025}; however, none of these approaches
address the 2D setting or provide systematic OOD evaluation.

Despite these advances, three critical gaps persist:
\begin{enumerate}
  \item \textbf{Parameter efficiency.}
    ALIGNN's 4.03\,M parameters are disproportionate for a dataset of
    ${\sim}10$\,k samples, risking overfitting and slow inference.
  \item \textbf{Physical inductive bias for defects.}
    Standard message-passing GNNs are bounded by a fixed cutoff radius
    $r_\mathrm{cut}$, yet point defects in 2D materials induce strain fields
    and charge-density perturbations extending well beyond typical
    $r_\mathrm{cut} = 5$--$8$\,\AA{}~\cite{fang2025}.
    Capturing such long-range effects requires either very deep networks
    (information dilution) or explicit global interactions.
  \item \textbf{Generalization assessment.}
    All prior work on IMP2D reports only random-split accuracy---a measure
    of interpolation, not extrapolation~\cite{li2025ood}.
    No study has evaluated whether these models transfer to \emph{unseen
    host materials} within the same chemical family, let alone across
    families.
\end{enumerate}

%----------------------------------------------------------------------
\subsection{Our contributions}
%----------------------------------------------------------------------

We address these gaps by introducing \dart{} (Defect-Aware Radial
Transformer), a compact model with explicit physical inductive biases for 2D
defect systems, accompanied by a rigorous multi-tier generalization
assessment.
Our contributions are:

\begin{enumerate}
  \item \textbf{\dart{} architecture} (\S\ref{sec:arch}):
    A 0.75\,M-parameter model combining local bond--angle message passing
    (three-body interactions within $r_\mathrm{cut}=5$\,\AA{}) with global
    multi-head self-attention biased by periodic minimum-image distances
    ($d_\mathrm{max}=12$\,\AA{}), augmented by an explicit defect-site
    embedding.

  \item \textbf{Systematic training pipeline} (\S\ref{sec:training}):
    MAE loss, cosine warmup, stochastic weight averaging (SWA), ct-UAE
    pretrained atomic embeddings~\cite{ctuae2025}, and multi-source data
    fusion (IMP2D + three auxiliary databases) yield a single-model MAE of
    0.407\eV{} and a 29-model ensemble of 0.359\eV{} ($-34\%$ vs.\ ALIGNN).

  \item \textbf{Graduated Generalization Assessment} (\S\ref{sec:ood}):
    A three-tier protocol---in-distribution, intra-family
    leave-one-host-out (LOHO), and compositional block-out---quantifying
    \dart{}'s deployment boundaries.
    Intra-family OOD degrades by only 1.5$\times$ (0.54\eV{}), far below
    the 4.9$\times$ naive-baseline collapse.

  \item \textbf{Calibrated uncertainty quantification}
    (\S\ref{sec:uq}):
    Temperature-scaled ensemble variance achieves 93.4\% coverage at the
    90\% nominal level.
\end{enumerate}
